\section*{الفصل \ref{ch:b1:radiation-heat-transport} --- الإشعاع ونقل الحرارة}
% القيم: calc/ch27.py

\begin{solution}{exo:b1:radiation-heat-transport:1}
(أ)~$S_0/\pi = 1361/\pi = \qty{433.2}{W.m^{-2}}$. (ب)~عند خطّ الاستواء في يوم اعتدال
$\bar Q = (S_0/\pi)(a/r)^2$، إذن $(a/r)^2 = 436.7/433.2 = 1.0080$ و$r/a =
1.0080^{-1/2} = 0.9960$: الأرض أقرب إلى الشمس ب\qty{0.4}{\percent} من بُعدها المتوسط. (ج)~$\bar Q = 433.2 \times 0.9683 \times \cos\ang{23.44} =
\qty{384.9}{W.m^{-2}}$.
\end{solution}

\begin{solution}{exo:b1:radiation-heat-transport:2}
(أ)~$415.6/172.4 = 2.41$. (ب)~تحصر المنطقة بين الموازيين $\pm\varphi$ مساحة $2\pi a_\oplus^2 \times 2\sin\varphi$، أي كسرًا قدره $\sin\varphi$ من الكرة:
$\sin\ang{37.0} = 0.602$، \qty{60}{\percent}. (ج)~على محور بدلالة خطّ العرض، تبدو النطاقات قرب القطبين عريضة مثل تلك قرب خطّ الاستواء، لكنها تغطّي مساحة أقلّ بكثير: فالمنطقة شمال \ang{60} ليست سوى $1 - \sin\ang{60} =
\qty{13.4}{\percent}$ من السطح. وبالنسبة إلى $\sin\varphi$، تملأ المدارات معظم المحور.
\end{solution}

\begin{solution}{exo:b1:radiation-heat-transport:3}
(أ)~$T_e = (\text{OLR}/\sigma)^{1/4}$: \qty{256.4}{K}، و\qty{241.3}{K} و
\qty{254.5}{K}. (ب)~الإشعاع الشمسي الممتص: $287.7/75.8 = 3.80$؛ OLR: $244.9/192.3 =
1.27$. (ج)~يتبع \omterm{def:b1:radiation-heat-transport:net-radiation}{الإشعاع الصافي} الإشعاع الشمسي الممتص: موجب وكبير في المدارات، وسالب بشدّة عند \omterm{def:b1:spherical-astronomy:geographic}{خطوط العرض} المرتفعة. فإن OLR، الذي يتغيّر قليلًا، لا يعوّض تناقص الإشعاع الشمسي الممتص.
\end{solution}

\begin{solution}{exo:b1:radiation-heat-transport:4}
(أ)~$2.82 \times 365.25 = \qty{1030}{mm}$، أي \qty{1.03}{m} في السنة. (ب)~$1.030 \times
78.1 = \qty{80.4}{W.m^{-2}}$.
(ج)~يبلغ $LE$ عند التحليل المُعاد \qty{2}{W.m^{-2}} أكثر. على مدى سنوات كثيرة لا يستطيع الغلاف الجوي أن يواصل تراكم بخار الماء: فكل ماء تبخّر يعود إلى الهبوط، لذا يجب أن يتساوى $LE$ و$LP$ في المتوسط العام. والفرق عدم توازن للنموذج، كعدم توازنه الإشعاعي.
\end{solution}

\begin{solution}{exo:b1:radiation-heat-transport:5}
(أ)~$\bar Q = S_0(a/r)^2\sin\delta = 1361 \times 0.9683 \times 0.3978 =
\qty{524.2}{W.m^{-2}}$. (ب)~$1361 \times 1.0334 \times 0.3978 =
\qty{559.5}{W.m^{-2}}$؛ والنسبة $1.0334/0.9683 = 1.067$. (ج)~للقطب ضوء شمس \num{24} ساعة من أصل \num{24} على ارتفاع ثابت، فمتوسط $\sin h_\odot$ عليه طوال اليوم هو $\sin\delta = 0.398$. وعند خطّ الاستواء تكون الشمس مرتفعة عند الظهر لكنها تغيب نصف الوقت؛ ومتوسط $\sin h_\odot$ على الساعات الأربع والعشرين هو $\cos\delta/\pi =
0.292$. والنسبة $\pi\tan\delta = 1.362$
(\cref{prop:b1:radiation-heat-transport:polar-summer}).
\end{solution}

\begin{solution}{exo:b1:radiation-heat-transport:6}
(أ)~$\cos h_0 = -\tan\ang{40}\tan\ang{23.44} = -0.3638$، $h_0 = \ang{111.33} =
\qty{1.943}{rad}$؛ ويدوم النهار $24 \times 111.33/180 = \qty{14.84}{h}$؛ و$\bar Q =
(1361/\pi) \times 0.9683 \times (1.943 \times 0.6428 \times 0.3978 + 0.7660 \times
0.9175 \times 0.9315) = \qty{483.0}{W.m^{-2}}$. (ب)~$\cos h_0 = +0.3638$، $h_0 =
\ang{68.67}$، و\qty{9.16}{h}؛ $\bar Q = \qty{155.9}{W.m^{-2}}$. (ج)~النسبة
$483.0/155.9 = 3.10$. والمسافة تُسهم ب$0.9683/1.0334 = 0.937$: وهي تعمل ضد الصيف الشمالي. والهندسة، أي \omterm{def:b1:spherical-astronomy:horizontal}{ارتفاع الشمس} وطول النهار، تُسهم ب$3.10/0.937 = 3.31$، أي كل التباين الفصلي وأكثر.
\end{solution}

\begin{solution}{exo:b1:radiation-heat-transport:7}
(أ)~$B = H/LE$: \num{0.093} عند خطّ الاستواء، و\num{0.374} عند \ang{25} شمالًا،
\num{0.338} عند \ang{45} شمالًا، و\num{-0.152} عند \ang{70} شمالًا. (ب)~\ang{25} شمالًا: المتوسط النطاقي يشمل الصحاري الحارّة الكبرى، أرضًا جافّة بقليل من الماء ليتبخّر، حيث تسخّن الأرض المسخّنة الهواء بدلًا من ذلك. (ج)~$H < 0$ تعني أن \omterm{def:b1:heat-budget:heat-flow}{تدفّقات الحرارة} تتّجه من الهواء إلى السطح: ففي المتوسط السنوي يكون الثلج والجليد البحري أبرد من الهواء فوقهما، فيُسخَّن الهواء بهواء يأتي من خطوط عرض أدنى.
\end{solution}

\begin{solution}{exo:b1:radiation-heat-transport:8}
(أ)~الإقرار: $\alpha = 100/340 = 0.294$، $R = 340 - 100 - 239 =
\qty{+1}{W.m^{-2}}$ (تقريب \num{0.7}). التحليل المُعاد: $\alpha = 117.0/341.9 =
0.342$، $R = 341.9 - 117.0 - 237.7 = \qty{-12.8}{W.m^{-2}}$ بهذه الحدود المقرّبة،
\qty{-12.9}{W.m^{-2}} بالحدود غير المقرّبة ل\cref{ex:b1:radiation-heat-transport:global}. (ب)~يمتصّ التحليل المُعاد $224.8$ مقابل $240$، أي \qty{15.2}{W.m^{-2}} أقلّ، لأنه يعكس \qty{17}{W.m^{-2}} أكثر؛ وOLR أدنى ب\qty{1.3}{W.m^{-2}} فقط. فالانعكاس الزائد يفسّر عدم التوازن. (ج)~تكامل $R$ من القطب الجنوبي إلى القطب الشمالي يعطي التكامل العام، $-12.9 \times 4\pi a_\oplus^2 = -12.9 \times
5.10\times10^{14} = \qty{-6.6}{PW}$. والنقل عبر القطب الشمالي مستحيل: فلا شيء وراءه، فالموازي نقطة. ولا تصحّ القضية إلا إذا كان $\int R\,\mathrm dA = 0$، وهو ما تقتضيه الحالة المستقرّة.
\end{solution}

\begin{solution}{exo:b1:radiation-heat-transport:9}
(أ)~المساحة هي $2\pi a_\oplus^2 \times 2\sin\ang{30} = 2\pi a_\oplus^2 =
\qty{2.55e14}{m^2}$، أي نصف الأرض؛ والفائض هو $33 \times 2.55\times10^{14} =
\qty{8.4}{PW}$. (ب)~\qty{4.2}{PW} عبر كل موازٍ. (ج)~المجموعان يساويان
\qty{8.5}{PW}، أي الفائض، لكنهما غير متساويين: فنصف الكرة الجنوبي يصدّر \qty{0.5}{PW} أكثر من الشمالي. ولا يُقتسم الفائض عند خطّ الاستواء: ففي هذا التحليل المُعاد، \qty{0.36}{PW} من فائض المدارات الشمالية يعبر خطّ الاستواء نحو الجنوب.
\end{solution}

\begin{solution}{exo:b1:radiation-heat-transport:10}
(أ)~اكتب المتوسط اليومي ك$\bar Q = (a/r)^2 F(\varphi, \delta)$، حيث $F$ القوس في \cref{thm:b1:radiation-heat-transport:daily} مضروبًا في $S_0/\pi$. خلال $\mathrm dt$ تكون الطاقة المستقبلة لكل متر مربّع هي $\bar Q\,\mathrm dt$، و$\mathrm dt = r^2\,\mathrm d\lambda/h$، إذن $\bar Q\,\mathrm dt = (a^2/h)
F(\varphi, \delta(\lambda))\,\mathrm d\lambda$: يحذف العامل $r^2$ نفسه. (ب)~على مدار السنة تكون الطاقة $(a^2/h)\int_0^{2\pi} F(\varphi,
\delta(\lambda))\,\mathrm d\lambda$. تبقى الصيغة دون تغيير حين يغيّر $\varphi$ و$\delta$ الإشارة معًا، $F(-\varphi, \delta) = F(\varphi, -\delta)$، و$\delta(\lambda + \pi) = -\delta(\lambda)$، لأن $\sin\delta = \sin\varepsilon
\sin\lambda$. وبإزاحة $\lambda$ بمقدار $\pi$ في التكامل يتبيّن أن $-\varphi$ يتلقّى الطاقة نفسها التي يتلقّاها $\varphi$. (ج)~يدوم الربيع والصيف الشماليان \num{7.6} يومًا أطول، لأن الأرض قرب الأوج وتتحرّك ببطء حينئذٍ؛ والصيف الجنوبي أشدّ إشراقًا لكنه أقصر. وبموجب (أ)، يلغي التأثيران أحدهما الآخر تمامًا.
\end{solution}

\begin{solution}{exo:b1:radiation-heat-transport:11}
(أ)~$\langle (a/r)^2\rangle = \frac1P\int_0^P \frac{a^2}{r^2}\,\mathrm dt =
\frac1P\int_0^{2\pi}\frac{a^2}{r^2}\frac{r^2}{h}\,\mathrm d\lambda = \frac{2\pi
a^2}{hP} = \frac{2\pi a^2}{2\pi a b} = \frac ab$، و$b = a\sqrt{1-e^2}$.
(ب)~المتوسط العام اللحظي $\frac14 S_0(a/r)^2$ يتوسّط إلى $\frac14 S_0
\langle(a/r)^2\rangle = S_0/(4\sqrt{1-e^2})$. (ج)~$340.25/\sqrt{1 - 0.0167^2} =
\qty{340.30}{W.m^{-2}}$، أي \qty{0.05}{W.m^{-2}} أو \qty{0.014}{\percent} فوق
$S_0/4$.
\end{solution}

\begin{solution}{exo:b1:radiation-heat-transport:12}
(أ)~الغلاف الجوي $0.78 \times 4.10 = \qty{3.20}{PW}$، والمحيط $0.22 \times 4.10 =
\qty{0.90}{PW}$. (ب)~يُوجد أن الغلاف الجوي وحده ينقل أكثر، \qty{5.0}{PW}، من الغلاف الجوي والمحيط معًا في التحليل المُعاد، \qty{4.02}{PW}: فالنقل الكلّي للتحليل المُعاد أقصر من ذلك بخُمس على الأقلّ ($1 - 4.02/5.0 = 0.20$). وإشعاعه الصافي يتغيّر أقلّ من اللازم بين المدارات والقطبين. (ج)~طرح ثابت من $R$ يجعل تكامله العام صفرًا، لكنه يترك تغيّره بخطّ العرض على حاله؛ والخطأ يكمن في ذلك التغيّر، في ألبيدو غيوم النموذج وأسطحه عند كل خطّ عرض، لا في المتوسط العام وحده.
\end{solution}

\begin{solution}{pb:b1:radiation-heat-transport:1}
\begin{enumerate}
  \item يونيو: $\cos h_0 = -0.3638$، $h_0 = \ang{111.33}$، $\bar Q =
        \qty{483.0}{W.m^{-2}}$. ديسمبر: $h_0 = \ang{68.67}$، $\bar Q =
        \qty{155.9}{W.m^{-2}}$. مارس: $h_0 = \ang{90}$، $\bar Q = (1361/\pi)
        \times 1.0081 \times \cos\ang{40} = \qty{334.6}{W.m^{-2}}$.
  \item يونيو: $\tan\ang{80}\tan\ang{23.44} = 2.46 > 1$، نهار قطبي، $h_0 = \pi$:
        $\bar Q = 1361 \times 0.9683 \times \sin\ang{80}\sin\ang{23.44} =
        \qty{516.3}{W.m^{-2}}$. ديسمبر: ليل قطبي، $\bar Q = 0$. مارس:
        $(1361/\pi) \times 1.0081 \times \cos\ang{80} = \qty{75.8}{W.m^{-2}}$.
  \item عند \ang{80} شمالًا لا تغيب الشمس ($h_0 = \ang{180}$، 24 ساعة) بينما عند
        \ang{40} شمالًا تُضيء \qty{14.8}{h}. وارتفاعها أدنى، بين
        \ang{13.4} و\ang{33.4}، لكن الحدّ $h_0\sin\varphi\sin\delta$، الذي ينمو بطول النهار وخطّ العرض، هو الذي يغلب.
  \item قيم النطاقات تتجاوز القيم المركزية ب\qty{0.5}{\percent}،
        \qty{0.8}{\percent}، \qty{1.7}{\percent}، \qty{1.7}{\percent} و
        \qty{1.3}{\percent}. أولًا، يستخدم التحليل المُعاد إشعاعًا شمسيًّا أكبر قليلًا (السؤال~5)، ب\qty{0.5}{\percent} أكثر. وثانيًا، متوسط النطاق ليس القيمة عند المركز: فنصف كل نطاق الأقرب إلى خطّ الاستواء له مساحة أكبر وضوء شمس أكثر من نصفه الأقرب إلى القطب.
  \item بموجب \cref{prop:b1:radiation-heat-transport:global}، $S_0 = 4\sqrt{1-e^2}
        \times 341.9 = \qty{1367}{W.m^{-2}}$، أي \qty{0.47}{\percent} أكثر من \qty{1361}{W.m^{-2}}.
  \item $A = 2\pi a_\oplus^2(1 - \sin\ang{40}) = 2\pi \times (6.371\times10^6)^2
        \times 0.3572 = \qty{9.110e13}{m^2}$، أي كسر $(1 - \sin\ang{40})/2 =
        0.179$ من الأرض.
  \item $(1 - \alpha)\bar Q$: \num{198.0}، \num{159.3}، \num{124.7}، \num{92.5}،
        \num{74.7} \unit{W.m^{-2}}.
  \item $R$: \num{-28.1}، \num{-51.0}، \num{-76.6}، \num{-104.9}، \num{-116.9}
        \unit{W.m^{-2}}. كل النطاقات الخمسة في عجز؛ والنطاق القطبي
        \ang{80}--\ang{90} له أكبر عجز.
  \item $T_e = (\text{OLR}/\sigma)^{1/4}$: \num{251.3}، \num{246.8}، \num{244.1}،
        \num{242.9}، \num{241.1} \unit{K}.
  \item الإشعاع الممتص: $198.0/74.7 = 2.65$؛ OLR: $226.1/191.6 = 1.18$.
  \item الشمس أدنى، والضوء الوارد بزاوية مماسّة ينعكس أكثر بفعل السحب والماء؛ والثلج والجليد البحري، وهي أسطح لامعة، يغطّيان مساحة أكبر من الأرض والبحر نحو القطب.
  \item يتناسب OLR مع $T_e^4$، و$T_e$ يتغيّر ب\qty{10}{K}، أي \qty{4}{\percent}، فقط عبر القلنسوة: فالعروض العالية، العروض العالية، تُبقيها الحرارة المجلوبة من الجنوب أدفأ بكثير مما يسمح به ضوء شمسها وحده.
  \item ب\unit{m^2}: \num{3.143e13}، \num{2.550e13}، \num{1.879e13}،
        \num{1.151e13}، \num{0.387e13}؛ ومجموعها \qty{9.110e13}{m^2}.
  \item $R_iA_i$: \num{-0.883}، \num{-1.300}، \num{-1.439}، \num{-1.207}،
        \num{-0.453} \unit{PW}؛ القلنسوة: \qty{-5.28}{PW}.
  \item $\alpha_\text{القلنسوة} = \sum\alpha_i\bar Q_iA_i/\sum\bar Q_iA_i = 0.400$،
        مقابل \num{0.342} للكرة كلّها.
  \item $-12.90 \times 5.10\times10^{14} = \qty{-6.58}{PW}$: الأرض ستفقد الحرارة بهذا المعدّل وتبرد. أمّا عدم التوازن المقدَّر فهو
        $+0.7 \times 5.10\times10^{14} = \qty{+0.36}{PW}$، أصغر بثماني عشرة مرّة
        وبإشارة معاكسة: عدم توازن التحليل المُعاد خطأ في النموذج.
  \item $R_i + 12.90$ المصحَّح: \num{-15.2}، \num{-38.1}، \num{-63.7}،
        \num{-92.0}، \num{-104.0} \unit{W.m^{-2}}؛ القدرات \num{-0.478}،
        \num{-0.972}، \num{-1.197}، \num{-1.059}، \num{-0.403} \unit{PW}؛ القلنسوة:
        \qty{-4.108}{PW}، بعجز متوسط قدره \qty{45.1}{W.m^{-2}}.
  \item لا تتغيّر طاقة القلنسوة من سنة إلى أخرى، فما تخسره عند أعلى الغلاف الجوي يجب أن يدخلها، ولا منفذ للدخول إلا عبر الموازي الأربعين: $T(\ang{40}) = -2\pi a_\oplus^2\int_{\ang{40}}^{\ang{90}}
        R\cos\varphi\,\mathrm d\varphi$.
  \item $T(\ang{40}\,\text{شمالًا}) = +\qty{4.11}{PW}$، نحو الشمال.
  \item $LP = 2.4655\times10^6 \times 1.95/86400 = \qty{55.6}{W.m^{-2}}$، أكثر من
        $LE = \qty{42.4}{W.m^{-2}}$: يهطل على القلنسوة ماء أكثر مما يتبخّر منها؛ والفرق يصل بخارًا.
  \item $(LE - LP) - 2.07 = 42.4 - 55.6 - 2.07 = \qty{-15.3}{W.m^{-2}}$؛ مضروبًا في
        مساحة القلنسوة: \qty{-1.40}{PW}. يجلب الغلاف الجوي \qty{1.40}{PW}
        من الحرارة الكامنة عبر \ang{40} شمالًا.
  \item $1.40/4.11 = 0.34$: ثلث.
  \item المحيط $0.22 \times 4.11 = \qty{0.90}{PW}$؛ الغلاف الجوي $0.78 \times 4.11 =
        \qty{3.20}{PW}$، منها \qty{1.40}{PW} حرارة كامنة و\qty{1.80}{PW}
        محمولة كحرارة محسوسة وطاقة كامنة للهواء.
  \item حدّ أدنى. فالمجموع الكلّي للتحليل المُعاد عند \ang{43} شمالًا، \qty{4.02}{PW}، أقلّ من الغلاف الجوي وحده في الدراسة، \qty{5.0}{PW}؛ والمجموع الحقيقي أكبر.
  \item نحو \qty{4.11}{PW}، أي \qty{4.11e15}{W}، وفق إشعاع التحليل المُعاد، ثلثه حرارة كامنة؛ والقيمة الحقيقية أكبر أيضًا.
\end{enumerate}
\end{solution}
